\begin{ericbox}{new section draft --- related literature ($\approx$750 words, RP requires it)}
To follow the introduction as Section 2, pushing everything else down one.
Citations are plain text here; on adoption they become \verb|\citep| calls
and the named entries return from
\texttt{\_src/bibliography\_parked.tex}. Edit freely --- this is the one
section referees will read as our claim about the field.

\medskip\noindent\textbf{Related literature}

\medskip\noindent\emph{Trademarks as innovation data.} A trademark filing
is a dated, costly, public disclosure made by the millions of firms that
never patent, and a literature has grown around reading it as an
innovation indicator (Castaldi 2020; Willeke, Block and Lambrecht 2026
review it). Trademark counts and breadth predict venture valuations
(Block, De Vries, Schumann and Sandner 2014); filing timing tracks
innovation processes (Seip, Castaldi, Flikkema and de Man 2018); and
trademarks mark the product-market side of innovation that patents miss
(Flikkema, Castaldi, de Man and Seip 2019; Flikkema, Castaldi and de Man
2025). Census-linked work reads registration as a correlate of firm
growth and survival (Dinlersoz, Goldschlag, Myers and Zolas 2018;
Dinlersoz, Goldschlag, Yorukoglu and Zolas 2023). Nearly all of it counts
filings, classes, or renewals; the descriptions' text enters, if at all,
as breadth or similarity. This paper reads the text itself, and reads it
against time.

\medskip\noindent\emph{Text measures of novelty.} Scoring a document by
its divergence from a reference corpus before and after it comes from
cultural analytics (Murdock, Allen and DeDeo 2017, on Darwin's notebooks;
Barron, Huang, Spang and DeDeo 2018, on the French Revolution's debates).
In economics and finance, text-based measures locate firms in product
space (Hoberg and Phillips 2016), date technological innovation in
patents over the long run (Kelly, Papanikolaou, Seru and Taddy 2021), and
track the vintage of patent language (Kalyani 2024). The construction
here differs in its object --- the composition of commercial offerings,
not ideas or claims --- and in its two-sided reference: every filing is
scored against its own industry's language both before and after its
date, which is what lets ``unusual'' be separated from ``early.'' The
separation matters because the two components predict opposite things.

\medskip\noindent\emph{Pioneering and optimal distinctiveness.} Whether
market pioneers are advantaged or burdened is an old, unresolved question
measured mostly on named markets and survivor-prone samples (Golder and
Tellis 1993; Semadeni and Anderson 2010 on the follower's dilemma in
services). A parallel literature finds evaluators reward conformity with
a penalty for extremity in both directions --- the optimal
distinctiveness of Zuckerman (1999), Hsu (2006), Zhao, Fisher, Lounsbury
and Miller (2017), and Taeuscher and Rothe (2021). This paper meets both
at census scale: every entrant's language is observed, whether or not the
product survived to be remembered, and the record supplies two different
evaluators --- an examiner at registration and the market at the
five-year proof of continued use --- who turn out to price the two
components differently. The examiner rewards middling conformity and
cannot observe lead; the market is indifferent to conformity once theme
importation is held, and penalizes having been early.

\medskip\noindent\emph{What this paper adds.} Three things. A measure:
atypicality and lead, computable for any filing from public data, with
its validity conditions stated (lead is robust across vocabulary
partitions, weightings and windows; atypicality's protective effect is a
property of global partitions). A set of facts: leading language predicts
cancellation for non-use, in three of four industries, twice as strongly
in the industries and eras where language moved fastest, with no
offsetting reward at funding, reporting, or listing. And two hazards for
the growing trademark-text literature: divergence measured on words
largely measures description length, and the registered description is a
conformed document, rewritten toward the class's typical language by the
applicants who try again --- so text indicators built on registered marks
inherit both.
\end{ericbox}
